\input{../tex-inputs/latex-leseansicht-vorspann}

\section[Arthur und Olga Schnitzler an Gustav Schwarzkopf, 2. 7. 1909]{L04520 Arthur und Olga Schnitzler an Gustav Schwarzkopf, 2. 7. 1909}
\nopagebreak\mylabel{L04520v}
\rehead{ }\normalsize\beginnumbering\briefempfaengerindex{Schwarzkopf, Gustav@\textsc{Schwarzkopf, Gustav}!zzzSchnitzler, Olga@\emph{von Olga Schnitzler}!1909-07-021@{2. 7. 1909}|(be}\briefempfaengerindex{Schwarzkopf, Gustav@\textsc{Schwarzkopf, Gustav}!zzzSchnitzler, Arthur@\emph{von Arthur Schnitzler}!1909-07-021@{2. 7. 1909}|(be}
\toendnotes[C]{\smallbreak\pagebreak[2]}
\correspDesc{Versand  durch Arthur Schnitzler, Olga Schnitzler am 2. 7. 1909 in Edlach
\newline{}Übermittlung  am 2. 7. 1909 in Edlach
\newline{}Erhalt  durch Gustav Schwarzkopf im Zeitraum [3. 7. 1909 – 7. 7. 1909?] in Wien}\toendnotes[C]{\smallbreak}
\Standort{DLA, A:Schnitzler, HS.1985.1.1897.}
\physDesc{Bildpostkarte, 409 Zeichen
\newline{}Handschrift Arthur Schnitzler: Bleistift, deutsche Kurrent
\newline{}Handschrift Olga Schnitzler: Bleistift, lateinische Kurrent
\newline{}Versand: Stempel: »\nobreak{}\oindex{Edlach@\textbf{Edlach}|pwk}Edlach b. Reichenau in N.Oe., \textcolor{gray}{2} 7 09, 2–6N\nobreak{}«.  }\toendnotes[C]{\smallbreak}
\pstart
           \noindent{}\centering{}{\pb}\textcolor{gray}{\textbf{Gruss aus Edlach\oindex{Edlach@\textbf{Edlach}|pw}}}\hspace*{2.5em}\textcolor{gray}{\textbf{Hôtel Edlacherhof}}\oindex{Hotel Edlacherhof@\textbf{Hotel Edlacherhof}|pw}\pend
           \pstart{}{\pb}Hrn \textsc{Gustav Schwarzkopf}\pend{}\pstart{}\textsc{Wien I.\oindex{I., Innere Stadt@\textbf{I., Innere Stadt}|pw}}\pend{}\pstart{}\textsc{Tiefer Graben 23}\oindex{Wien@\textbf{Wien}!I., Innere Stadt@\textbf{I., Innere Stadt}!Tiefer Graben 23@\textbf{Tiefer Graben 23}|pw}\pend{}{\bigskip}\vspace{1em}
\pstart
           \noindent{}Lieber Guſtav, im Hotel Rax\oindex{Hotel Rax@\textbf{Hotel Rax}|pw} gibt es ſehr anſtändige \uline{Manſarden zu 2 Kronen},
              für die man nur rechtzeitig optiren
              müßte. Ich nannte dem Kellner\pwindex{?? [Kellner im Hotel Rax] @\textsc{?? [Kellner im Hotel Rax]}|pwv} Ihren
              Namen – er{ }ſagte: Aha, Bilanz der Ehe\pwindex{Schwarzkopf, Gustav 7.\,11.\,1853 Wien – 13.\,11.\,1939 ebd.@\textsc{Schwarzkopf, Gustav} (7.\,11.\,1853 Wien – 13.\,11.\,1939 ebd.)!Bilanz der Ehe. Novellistische Studien. 2 Bde.@\strich\emph{Die Bilanz der Ehe. Novellistische Studien. 2 Bde.}|pw}, –
              und erwartet Ihre näheren Befehle!
              Ich auch. Herzlichſt Ihr\pend
           \pstart \spacefill\mbox{A.}\pend{}
\pstart
           2.\,7.\,09.\pend
           \selectlanguage{ngerman}\vspace{1em}\pstart {[}hs. O. Schnitzler:{]} \spacefill\mbox{OlgaS.}\pend{}
\pstart
           \noindent{}{\pb}Der Kellner\pwindex{?? [Kellner im Hotel Rax] @\textsc{?? [Kellner im Hotel Rax]}|pwv} sagte in Wirklichkeit: »Den Namen
      kann man sich scho merken, – is a berühmter
      Name aus Deutschland\oindex{Deutschland@\textbf{Deutschland}|pw}!«\pend
           \selectlanguage{ngerman}\endnumbering\briefempfaengerindex{Schwarzkopf, Gustav@\textsc{Schwarzkopf, Gustav}!zzzSchnitzler, Olga@\emph{von Olga Schnitzler}!1909-07-021@{2. 7. 1909}|)be}\briefempfaengerindex{Schwarzkopf, Gustav@\textsc{Schwarzkopf, Gustav}!zzzSchnitzler, Arthur@\emph{von Arthur Schnitzler}!1909-07-021@{2. 7. 1909}|)be}\mylabel{L04520h}
\begin{anhang}
\end{anhang}\newcommand{\dateiname}{L04520}\newcommand{\titel}{Arthur und Olga Schnitzler an Gustav Schwarzkopf, 2. 7. 1909}\newcommand{\editorInnen}{Herausgegeben von Jahnke, SelmaMüller, Martin Anton}\input{../tex-inputs/latex-leseansicht-abspann}
